%% ===========================================================================
%% ICIEV 2026 regular-paper submission (revised draft, v5 final candidate)
%%
%% Template: accepted ICIEV 2026 papers are published in IJCVSP, so the
%% official IJCVSP two-column class is used (ijcvsplatex.cls, unmodified).
%% Regular papers: 6--10 pages. Review is double-blind.
%%
%% COMPILER: XeLaTeX (Overleaf: Menu -> Settings -> Compiler -> XeLaTeX).
%% XeLaTeX is only needed to typeset the Bangla examples with the bundled Noto
%% Serif Bengali font; under pdfLaTeX Bangla strings become "[Bangla]".
%%
%% Header policy for the initial (double-blind) submission. Neither the ICIEV
%% submission page nor the IJCVSP author guide says how the template's
%% publication metadata should look at submission, so we keep everything the
%% official template prints that is true for an unpublished manuscript (IJCVSP
%% logo, ISSN, journal URL, "ORIGINAL ARTICLE", the template's running heads
%% "CNSER" / "Int. J. Computer Vision Signal Process.", and the "Contribution
%% of the Paper" line that the template's sample asks for), and hide only the
%% placeholder values that would be false before acceptance: the article
%% history dates, the volume/page range and the copyright line.
%% \reviewversiontrue does this; set \reviewversionfalse for camera-ready, when
%% the editors supply the real values.
%% ===========================================================================
\documentclass[doublecolumn]{ijcvsplatex}
\usepackage{graphicx}
\usepackage{amsmath}
\usepackage{amssymb}
\usepackage{booktabs}
\usepackage{multirow}
\usepackage{array}
\usepackage{xcolor}
\usepackage{etoolbox}
\usepackage{pgfplots}
\pgfplotsset{compat=1.17}
\usetikzlibrary{arrows.meta,fit,backgrounds,calc,patterns}
\usepackage{iftex}
\interdisplaylinepenalty=2500
\setlength{\emergencystretch}{1.5em}

%% ---- Review copy: hide publication-only metadata --------------------------
\newif\ifreviewversion
\reviewversiontrue
\makeatletter
\ifreviewversion
  \def\@received{}\def\@revised{}\def\@accepted{}\def\@online{}
  \patchcmd{\MaketitleBox}{Article History:}{}{}{}
  \patchcmd{\MaketitleBox}{\textit{$\copyright$ 2026, IJCVSP, CNSER. All Rights Reserved}}{}{}{}
  %% First-page head as in the template, without the placeholder volume/pages.
  \def\ps@pprintTitle{%
    \def\@oddhead{\footnotesize International Journal of Computer Vision and Signal Processing\hfill\@original}%
    \let\@evenhead\@oddhead
    \let\@oddfoot\@empty\let\@evenfoot\@empty}
  %% The template's even-page head carries the placeholder volume "16(1)";
  %% use its odd-page head (journal name only) on every page.
  \fancyhead[RO,RE]{Int. J. Computer Vision Signal Process.}
  \fancyhead[LO,LE]{CNSER}
\fi
%% fancyhdr's head width is fixed before the class sets the page geometry, so
%% the running head stopped short of the right margin; span the full text width.
\AtBeginDocument{\setlength{\headwidth}{\textwidth}}
\makeatother

%% ---- Bangla text ----------------------------------------------------------
\ifXeTeX
  \usepackage{fontspec}
  \newfontfamily\bengalifont[Path=fonts/,Script=Bengali,Scale=0.95]{NotoSerifBengali-Regular.ttf}
  \newcommand{\bn}[1]{{\bengalifont #1}}
\else
  \newcommand{\bn}[1]{\textit{[Bangla]}}
\fi

%% ---- Shorthands -----------------------------------------------------------
\newcommand{\qwen}{Qwen2-VL}
\newcommand{\medvlm}{MedVLM-R1}
\newcommand{\gemma}{Gemma~3}
\newcommand{\medgemma}{MedGemma}

%% Model colours: light = general model, dark = its medical version
\definecolor{cQwen}{RGB}{158,202,225}
\definecolor{cMedVLM}{RGB}{33,113,181}
\definecolor{cGemma}{RGB}{253,174,107}
\definecolor{cMedGemma}{RGB}{217,72,1}
\pgfplotsset{modelbars/.style={
  width=\columnwidth, height=4.6cm, ybar, bar width=4.2pt, ymin=0,
  xtick=data, yticklabel style={font=\footnotesize},
  ylabel style={font=\footnotesize},
  xticklabel style={font=\scriptsize, align=center},
  legend style={font=\scriptsize, at={(0.5,1.02)}, anchor=south, legend columns=4, draw=none},
  legend image code/.code={\draw[#1] (0cm,-0.08cm) rectangle (0.25cm,0.12cm);},
  enlarge x limits=0.12}}
%% The Gemma pair is hatched so that the two pairs stay distinguishable when
%% printed in black and white (light and dark bars have similar grey levels
%% across the two pairs).
\pgfplotsset{gemmapair/.style={postaction={pattern=north east lines, pattern color=black!55}}}

%% Clickable citations and cross-references, no colour, no boxes
\makeatletter
\@ifpackageloaded{hyperref}{}{\usepackage{hyperref}}
\makeatother
\hypersetup{hidelinks, pdfauthor={}}

\begin{document}

\begin{frontmatter}

\title{Does Medical Adaptation Help Small Vision-Language Models? A Paired Zero-Shot Study}

%% Double-blind review: keep anonymous until camera-ready.
\author{Anonymous Authors}
\address{Affiliation withheld for double-blind review}

\begin{abstract}
Medical vision-language models (VLMs) are usually built by further training a general-purpose VLM, on the assumption that this makes them better at medical images. We test this directly with two small open model pairs: \medvlm{} versus its base \qwen{}-2B, and \medgemma{}~4B versus its base \gemma{}~4B. All four models answer the same items, zero-shot, with the same prompts, on three tasks: benign/malignant classification of breast ultrasound (BUS-BRA, 400 images), closed-ended CT/MRI question answering (RadImageNet-VQA, 500 questions) and open-ended medical question answering in Bangla (BanglaMedVQA, 500 questions). Medical adaptation gave no consistent gain. \medvlm{} answered almost exactly like its base model on all three tasks. \medgemma{} improved clearly only in Bangla (+7.4 LAVE points), was not significantly better on breast ultrasound, and was slightly worse on RadImageNet-VQA. All models recognised modality and anatomy far better than disease, and often fell back on fixed answers: one label, “yes” or “no”, or one option letter.
\\
\textbf{Contribution of the Paper:} A paired comparison of medical VLMs against their own base models on three tasks, showing that medical adaptation does not consistently improve zero-shot performance and that answer biases and cross-lingual scoring can hide this behind headline accuracy.
\end{abstract}

\begin{keyword}
medical vision-language models \sep zero-shot evaluation \sep medical visual question answering \sep breast ultrasound \sep Bangla \sep answer bias
\end{keyword}

\end{frontmatter}

%% ===========================================================================
\section{INTRODUCTION}
\label{sec_intro}

Most medical vision-language models (VLMs) are not trained from scratch. They start from a general-purpose VLM and are trained further on medical images and text, through instruction tuning~\cite{li2023llavamed,chen2024huatuogptvision}, large-scale medical training~\cite{sellergren2025medgemma} or reinforcement learning on a small set of medical questions~\cite{pan2025medvlmr1}. The implicit claim is that the medical model is better at medical tasks than the model it started from.

That claim is rarely tested under matched conditions. New medical models are usually compared with other medical models or with larger proprietary systems, and when the base model is reported, its numbers often come from different prompts or evaluation code. Where the comparison has been controlled, the results are sobering: Jeong et al.~\cite{jeong2025limited} found that the medical VLMs they tested did not consistently beat their base models once each model was given its best prompt and statistical uncertainty was taken into account, and broad benchmarks have found general-purpose VLMs matching or beating medical ones~\cite{hu2024omnimedvqa,liu2025howfar}. Those studies covered earlier medical VLMs, English-only tasks and aggregate accuracy.

We ask the same question for two newer, small, open models that a group with limited resources could run locally: \emph{used zero-shot, does a medical VLM answer medical image questions better than its own base model?} \medvlm{}~\cite{pan2025medvlmr1} was built from \qwen{}-2B~\cite{wang2024qwen2vl} with reinforcement learning on 600 questions, and \medgemma{}~4B~\cite{sellergren2025medgemma} was built from \gemma{}~4B~\cite{gemma3} with large-scale medical training, so the two pairs represent very different amounts of adaptation. Within each pair the architecture is the same, so the comparison shows what the released medical weights change, although it cannot separate the ingredients of each adaptation recipe.

We use three tasks, each adding something the previous one cannot show:
\begin{itemize}
\item \textbf{Breast ultrasound classification} on the BUS-BRA dataset~\cite{gomezflores2024busbra}: one binary diagnostic decision with a known chance level. It shows whether a model uses the image at all.
\item \textbf{Closed-ended CT/MRI VQA} on RadImageNet-VQA~\cite{butsanets2025radimagenetvqa}: yes/no and multiple-choice questions that separate recognising anatomy from recognising disease, scored exactly.
\item \textbf{Open-ended Bangla VQA} on the BanglaMedVQA dataset~\cite{ahmed2026banglamedvqa}: free-text answers to questions in a low-resource language, where both the answers and their scoring are less constrained.
\end{itemize}

Neither medical model was consistently better than its base. \medvlm{} behaved almost as a copy of \qwen{}; \medgemma{} was clearly better than \gemma{} only in Bangla, and only on general questions. Looking at the answers rather than the scores explains much of this: several models gave nearly the same answer to every image, preferences for ``yes'', ``no'' or one option letter decided many disease questions, and the Bangla metric mis-scored correct answers given in English.

The main contributions of this paper are as follows:
\begin{enumerate}
    \item a paired zero-shot comparison of two recent medical VLMs with their own base models on three tasks, with identical items, prompts and decoding, and with all six differences tested and corrected for multiple comparisons
    \item an analysis of how the models fail, including constant answers, yes/no and option-letter preferences, and the gap between anatomy and disease questions
    \item an audit of the Bangla evaluation showing that an LLM-judged metric~\cite{manas2024lave} gives inconsistent credit to correct answers in the wrong language.
\end{enumerate}


%% ===========================================================================
\section{RELATED WORKS}
\label{sec_rl}

\subsection{Medical VLMs and their base models}
\textbf{Base models.} \qwen{}~\cite{wang2024qwen2vl} combines a vision transformer with the Qwen2 language model. It processes images at their native resolution, turning each image into a variable number of visual tokens, and was released at 2B, 8B and 72B parameters. \gemma{}~\cite{gemma3} is a family of open models from 1B to 27B parameters; the report introduces vision input through a SigLIP-based encoder, a longer context and improved multilingual ability over Gemma~2. We use the instruction-tuned 2B and 4B models.

%% Fig. 1 is defined here (page 2) so that it floats to the top of page 3,
%% next to its first mention in Section 3.1.
\begin{figure*}[tb]
\centering
\input{figures/pipeline_v5.tex}
\caption{Study design. (a) Each task uses a fixed, seeded sample that is identical for all four models. Images: BUS-BRA~\cite{gomezflores2024busbra} bus\_0007-l and bus\_0001-l (CC~BY~4.0) and ChestX-ray8~\cite{wang2017chestxray8} 00001170\_046 via BanglaMedVQA; RadImageNet-VQA images are not shown because its research-use agreement does not allow reproduction. (b) Arrows mark the adaptation done by each medical model's authors before this study; we used all four released models without further training. (c) Real answers from our result files to one item per task, for the images in (a); the RadImageNet-VQA item is a yes/no question, shown as its published template~\cite{butsanets2025radimagenetvqa} with the anatomy withheld. The paired comparisons and answer analyses run on these saved answers (Section~\ref{sec_stats}).}
\label{fig_pipeline}
\end{figure*}

\textbf{Medical VLMs.} Early medical VLMs adapted general models with biomedical image--text data: LLaVA-Med~\cite{li2023llavamed} fine-tuned LLaVA on PubMed Central figure--caption pairs and GPT-4-generated instructions, Med-PaLM~M~\cite{tu2023medpalmm} trained one large model on 14 biomedical tasks, and HuatuoGPT-Vision~\cite{chen2024huatuogptvision} was trained on 1.3 million medical VQA samples refined with GPT-4V. The two medical models we test represent two very different strategies for small models. \textbf{\medgemma{}}~\cite{sellergren2025medgemma} builds on \gemma{} with a medically tuned vision encoder (MedSigLIP) and further medical image and text training; its authors report gains of 2.6--10\% over the base models on out-of-distribution medical image question answering, measured by their own zero-shot evaluation. \textbf{\medvlm{}}~\cite{pan2025medvlmr1} applies Group Relative Policy Optimization~\cite{shao2024deepseekmath} to \qwen{}-2B using only 600 MRI multiple-choice questions. The reward checks that the model writes its reasoning in \texttt{<think>} tags and its final letter in \texttt{<answer>} tags, and that the letter is correct. The authors report an accuracy increase from 55.11\% to 78.22\% on MRI, CT and X-ray test questions.

\subsection{Evaluations of medical VLMs}
VQA-RAD~\cite{lau2018vqarad} and SLAKE~\cite{liu2021slake} are small radiology VQA sets often used for fine-tuning. Larger benchmarks have exposed the limits of current models. OmniMedVQA~\cite{hu2024omnimedvqa} collected questions from 73 datasets and found medical-specialised VLMs performing worse than general ones; GMAI-MMBench~\cite{chen2024gmaimmbench} evaluated 50 VLMs, the best of which (GPT-4o) reached 53.96\%; and a 2025 study of open models from 3B to 72B parameters found large general-purpose models matching or beating medical ones on several benchmarks, with reasoning consistently weaker than understanding~\cite{liu2025howfar}. Zhong et al.~\cite{zhong2025generalist} found that specialist medical VLMs keep an advantage when the imaging modality matches their training, but that efficiently fine-tuned generalists match or exceed them on most tasks. ProbMed~\cite{yan2024probmed} added negated questions with hallucinated findings and found strong models performing worse than random on specialised diagnostic questions, and a case study of GPT-4V~\cite{wu2023gpt4vmed} found good recognition of modality and anatomy but poor diagnosis. RadImageNet-VQA~\cite{butsanets2025radimagenetvqa}, one of our datasets, reported that VLMs struggle with fine-grained pathology and fall to near-random performance without the image.

These benchmarks compare many models, but seldom a medical model with its own base under identical conditions. The closest work to ours is Jeong et al.~\cite{jeong2025limited}, who compared medical language and vision-language models head-to-head with their bases, optimised the prompt separately for each model and accounted for statistical uncertainty. All medical VLMs they tested failed to improve consistently over their bases in zero- and few-shot settings. Their VLM evidence, however, covers two earlier models (LLaVA-Med and Med-Flamingo) on English VQA scored by exact-match accuracy. We apply the same paired logic to two newer models built with contrasting adaptation strategies, add a classification task with balanced labels and a low-resource language with free-text answers, and examine the answers themselves rather than only the scores. Where Jeong et al.\ gave each model its best prompt, we give both members of a pair the same prompt; the two designs answer complementary questions (Section~\ref{sec_design}).

\subsection{Multilingual medical VQA and answer biases}
Medical VQA has been evaluated almost only in English. IndoRad-VQA~\cite{indoradvqa2026} translated VQA-RAD into Indonesian and found gaps of 8--25\% between English and Indonesian, including answers in the wrong language. BanglaMedVQA~\cite{ahmed2026banglamedvqa}, the first Bangla medical VQA dataset, evaluated nine VLMs and found that medical VLMs (Med-LLaVA and MedGemma) underperformed general-purpose ones in zero-shot Bangla. Open-ended answers in such languages are often scored by an LLM judge, yet LLM judges agree poorly across languages, especially low-resource ones (average Fleiss' $\kappa\approx0.3$ over 25 languages~\cite{fu2025multilingualjudge}). Two further known biases matter for our analysis: VQA models can exploit answer priors such as favouring ``yes''~\cite{goyal2017vqav2}, and language models prefer certain option letters in multiple-choice questions~\cite{zheng2024mcqbias}.

%% ===========================================================================
\section{METHODOLOGY}
\label{sec_method}

\subsection{Study design}
\label{sec_design}
The study has one controlled comparison: each medical VLM against its own base model (Fig.~\ref{fig_pipeline}, Table~\ref{tab_models}). For each task, all four models receive the same items in the same order, with the same prompt and the same deterministic decoding; only model-specific loading and image preprocessing differ (Section~\ref{sec_inference}). No model is fine-tuned and the prompts contain no examples.

We chose zero-shot, direct-answer prompting for two reasons. It is how a team without labelled data or fine-tuning resources would use these models, and with the prompt held fixed, any difference within a pair comes from the weights rather than from prompt engineering. The alternative, tuning a prompt for each model as Jeong et al.~\cite{jeong2025limited} did, asks a different question: how good each model can be made. Our design therefore estimates what the medical weights change under one reasonable shared prompt, not the best each model can do; in particular, it does not use \medvlm{}'s reasoning format (Section~\ref{sec_limitations}).

\begin{table}[tb]
\centering
\caption{Model pairs. Each medical model was derived by its original authors from the general model in the same row (Hugging Face IDs in typewriter font). We used all four as released.}
\label{tab_models}
\footnotesize
\setlength{\tabcolsep}{3pt}
\begin{tabular}{@{}lll@{}}
\toprule
 & \textbf{General (base)} & \textbf{Medical version} \\
\midrule
$\sim$2B & \qwen{}-2B-Instruct~\cite{wang2024qwen2vl} & \medvlm{}~\cite{pan2025medvlmr1} \\
 & \scriptsize\texttt{Qwen/Qwen2-VL-2B-Instruct} & \scriptsize\texttt{JZPeterPan/MedVLM-R1} \\
 & & \scriptsize RL (GRPO) on 600 MRI questions \\
\midrule
$\sim$4B & \gemma{}~4B-it~\cite{gemma3} & \medgemma{}~4B-it~\cite{sellergren2025medgemma} \\
 & \scriptsize\texttt{google/gemma-3-4b-it} & \scriptsize\texttt{google/medgemma-4b-it} \\
 & & \scriptsize medical image and text training \\
\bottomrule
\end{tabular}
\end{table}

\subsection{Datasets and sampling}
\label{sec_data}
Fig.~\ref{fig_pipeline}(a) summarises the three samples and shows example images. All samples were drawn with seed 42 and saved, and the saved sample files are identical across the four runs of each task, so every model saw the same items. Every binary label set is balanced (benign/malignant, ``yes''/``no'' within each RadImageNet-VQA yes/no group, and 25 ``yes'' and 25 ``no'' BanglaMedVQA abnormality questions per source), so a model that always gives the same answer scores exactly chance and any answer preference becomes visible.

\textbf{BUS-BRA}~\cite{gomezflores2024busbra} contains 1,875 breast ultrasound images from 1,064 patients with biopsy-proven benign or malignant labels (1,268 benign, 607 malignant). We chose it because the task is a clinically meaningful binary decision with a clear chance level, the labels are biopsy-proven, the dataset is openly licensed (CC~BY~4.0), and case identifiers let us take one image per case. We sampled 400 images from 400 different cases, excluding 20 cases used in an earlier pilot run and 13 cases flagged during data checks. The final ID list was frozen and verified with a SHA-256 hash before each run.

\textbf{RadImageNet-VQA}~\cite{butsanets2025radimagenetvqa} is built from the CT and MRI part of RadImageNet~\cite{mei2022radimagenet}. Its benchmark split has 9,000 questions on 1,000 images from eight anatomical regions. We kept the 7,000 yes/no and multiple-choice questions, so that every answer could be scored exactly, and grouped them by format and topic into five groups (Table~\ref{tab_prompts}): anatomy and pathology multiple choice, and anatomy, abnormality and pathology yes/no. We drew 100 questions per group.

\textbf{BanglaMedVQA}~\cite{ahmed2026banglamedvqa} has 7,000 Bangla questions, half on chest X-rays from ChestX-ray8~\cite{wang2017chestxray8} and half on figures from biomedical papers collected in MedICaT~\cite{subramanian2020medicat}. The questions fall into five categories: \emph{modality}, \emph{organ}, \emph{abnormality} (yes/no), \emph{condition} and \emph{position}. The question--answer pairs were generated from image metadata with Gemini and translated into Bangla; two medical specialists accepted 97.2\% and 97.6\% of a 1,000-pair subset~\cite{ahmed2026banglamedvqa}. We drew 50 questions for each source--category cell from the training splits (3,500 questions per source), at most one per image within a cell. The original paper reports results on the test splits (200 questions per source), so our numbers are not directly comparable with theirs. Because we only run inference, using the train split does not involve any training; none of the four models is reported to have been trained on BanglaMedVQA.

\subsection{Prompts and question types}
Each task used one fixed prompt for all models (Table~\ref{tab_prompts}), which was not changed once evaluation started. For BanglaMedVQA we used the dataset authors' published zero-shot system prompt verbatim and passed the Bangla question as the user message, to stay as close as possible to the original benchmark. Table~\ref{tab_prompts} also describes each of the ten question types. The RadImageNet-VQA groups separate \emph{where} the image is from (anatomy) from \emph{what is wrong} (abnormality, pathology); BanglaMedVQA makes a similar split between \emph{general} (modality, organ) and \emph{specialised} (abnormality, condition, position) categories~\cite{ahmed2026banglamedvqa}.

\begin{table*}[tb]
\centering
\caption{Top: the prompts (verbatim; line breaks shown as $\hookleftarrow$). The image is placed before the text in one user turn. Bottom: the question types. RadImageNet-VQA questions are given as the templates published in the dataset paper~\cite{butsanets2025radimagenetvqa} (its Table~8), not as items from our sample, because the dataset's research-use agreement forbids reproducing any portion of it. BanglaMedVQA rows are real questions from our sample; English glosses for BanglaMedVQA are the dataset's own English versions; the models saw only the Bangla.}
\label{tab_prompts}
\footnotesize
\begin{tabular}{@{}>{\raggedright\arraybackslash}p{2.7cm}p{15.15cm}@{}}
\toprule
\multicolumn{2}{@{}l}{\textbf{Prompts}} \\
BUS-BRA & \texttt{You are a medical assistant. Examine this breast ultrasound image.$\hookleftarrow$ Classify the lesion as exactly ONE of the following:$\hookleftarrow$ - benign: non-cancerous lesion$\hookleftarrow$ - malignant: cancerous lesion$\hookleftarrow\hookleftarrow$ Respond with ONLY the label: benign or malignant.} \\
RadImageNet-VQA, MCQ & \texttt{You are a medical assistant. Look at this image and answer the following multiple-choice question with ONLY the letter of the correct option (e.g. 'B'), nothing else.$\hookleftarrow\hookleftarrow$ Question: \{question\}$\hookleftarrow$ Options:$\hookleftarrow$ A. \{A\}$\hookleftarrow$ B. \{B\}$\hookleftarrow$ C. \{C\}$\hookleftarrow$ D. \{D\}} \\
RadImageNet-VQA, yes/no & \texttt{You are a medical assistant. Look at this image and answer the following question with ONLY the single word 'Yes' or 'No'.$\hookleftarrow\hookleftarrow$ Question: \{question\}} \\
BanglaMedVQA \newline (system prompt of~\cite{ahmed2026banglamedvqa}; user turn = Bangla question) & \texttt{You are a careful, clinically grounded medical vision-language assistant. Task: You will be given a medical image and a single question about that image (e.g., modality, organ, abnormality, condition, or the position of the condition). Your job is to look at the image and provide the exact, minimal answer to the question. CRITICAL: You must respond with ONLY the exact answer. No explanations, no sentences, no extra words. STRICT OUTPUT RULES: - ONE word or short phrase only - NO sentences or explanations - NO "This is..." or "The image shows..." - NO punctuation unless part of the answer} \\
\midrule
\end{tabular}
\begin{tabular}{@{}>{\raggedright\arraybackslash}p{2.7cm}p{9.0cm}>{\raggedright\arraybackslash}p{4.3cm}c@{}}
\multicolumn{2}{@{}l}{\textbf{Example questions}} & \textbf{Reference answer} & \textbf{Chance} \\
\multicolumn{4}{@{}l}{\textit{RadImageNet-VQA (CT/MRI), 100 questions per group}} \\
Anatomy, MCQ & Which of the following anatomical regions is depicted in this scan? & one letter (A--D) & 25\% \\
Anatomy, yes/no & Is the \{anatomy\} visible in the image? & yes or no & 50\% \\
Abnormality, yes/no & Does the \{anatomy\} look abnormal? & yes or no & 50\% \\
Pathology, yes/no & Is the \{anatomy\} affected by \{pathology\}? & yes or no & 50\% \\
Pathology, MCQ & What is the most likely pathology in this image? & one letter (A--D) & 25\% \\
\multicolumn{4}{@{}l}{\textit{BanglaMedVQA (Bangla), 50 questions per category and source}} \\
Modality & \bn{এই গবেষণায় কোন ইমেজিং পদ্ধতি ব্যবহার করা হয়েছে?} (What imaging modality is used in this study?) & \bn{এক্স-রে} (X-ray) & -- \\
Organ & \bn{ছবিতে কোন অঙ্গটি দেখানো হয়েছে?} (What organ is depicted in the image?) & \bn{মস্তিষ্ক} (brain) & -- \\
Abnormality & \bn{কোনও অস্বাভাবিকতা আছে কি?} (Are any abnormalities present?) & \bn{হ্যাঁ} (yes) & 50\% \\
Condition & \bn{ছবিতে কী কী অবস্থা দেখা যাচ্ছে?} (What conditions are observed in the image?) & \bn{দুটি পাথর, এন্ডোক্লিপ} (two stones, endoclip) & -- \\
Position & \bn{নডিউলটি কোথায় অবস্থিত?} (Where is the nodule located?) & \bn{ডানদিকের মাঝখানে} (middle right) & -- \\
\bottomrule
\end{tabular}
\end{table*}

\subsection{Inference and answer parsing}
\label{sec_inference}
All models ran in bfloat16 with greedy decoding (no sampling), batch size~1 and seed~42, using each model's own chat template. Greedy decoding makes outputs deterministic for a given setup, and batch size~1 avoids padding effects. The output limit was 32 new tokens for BUS-BRA, 16 for RadImageNet-VQA and 96 for BanglaMedVQA (256 for \medvlm{}, so that a reasoning block, if produced, would not cut off the answer). The BUS-BRA runs pinned model revisions and library versions and ran on Kaggle T4 GPUs. In the two VQA runs, \medvlm{}'s own processor did not load with the installed \texttt{transformers} version, so we used the processor of its base model, \qwen{}-2B-Instruct; for BanglaMedVQA the Qwen processors were limited to $256$--$1280$ visual patches of $28\times28$ pixels.

Answers that could not be parsed would count as wrong, but every output was parseable. The BUS-BRA parser accepted only the bare label \texttt{benign} or \texttt{malignant} (case-insensitive, optionally inside \texttt{<answer>} tags); RadImageNet-VQA took the first standalone letter A--D or the first ``yes''/``no''; BanglaMedVQA answers were only stripped of whitespace, quotes and a final full stop or Bangla \emph{dari}. All 3,600 classification and closed-VQA outputs were a single word or letter, and no model produced \texttt{<think>} or \texttt{<answer>} tags in any task.

%% ===== REVISION (blue): blank-image control, Reviewers 9 and 11 ==============
\subsection{\textcolor{blue}{Blank-image control}}
\label{sec_blank_method}
{\color{blue}
To measure how much each answer depends on the image, we repeated all three tasks with the image replaced by a constant blank image and everything else unchanged: the same four models, items, prompts, decoding, answer parsing and metrics, and no retraining or tuning. For RadImageNet-VQA and BanglaMedVQA the input was a $512\times512$ solid white RGB image, passed in place of the real image for each of the 500 questions. For BUS-BRA we used both a white and a black $512\times512$ image. Because the input is identical for all 400 BUS-BRA items and decoding is deterministic, we ran five identical inferences per colour and model, observed the same output every time, and assigned that output to all 400 items; the BUS-BRA blank-image scores are therefore derived from five inferences per condition, not from 400 independent ones.

We compare each model's real-image and blank-image answers on the same items. Accuracy differences use the exact McNemar test and LAVE differences the paired permutation test of Section~\ref{sec_stats}, with a 95\% paired bootstrap interval. To ask whether a medical model uses the image more than its base model, we also compare their real-minus-blank LAVE differences item by item (bootstrap and permutation test as above). These checks are descriptive and are not among the six main hypotheses, so they are not Holm-adjusted. We did not run a shuffled-image control (real images paired with the wrong questions).
}
%% ===== end of revision ======================================================

\subsection{Metrics and statistics}
\label{sec_stats}
For BUS-BRA and RadImageNet-VQA we report accuracy with 95\% Wilson intervals~\cite{wilson1927} and a one-sided exact binomial test against chance. For BUS-BRA we add sensitivity and specificity, and for yes/no questions the share of ``yes'' answers, since a balanced set makes any answer preference visible.

For BanglaMedVQA we follow the original paper and report exact match (EM), BERTScore F1~\cite{zhang2020bertscore} (cased multilingual BERT-base, layer~9, no rescaling) and LAVE~\cite{manas2024lave}, in which an LLM judge rates each answer between 0 and 1 against the reference. As in the original paper, the judge was GPT-4.1-mini with the paper's judge prompt, which shows only the reference and the answer, not the question; we used one call per answer at temperature~0, and no call failed. LAVE is our main Bangla metric because EM is close to zero for all models. We also checked LAVE with two rule-based audits that use no judge: whether a yes/no answer is correct in either language, and whether the modality named in a modality answer is correct.

Each medical model is compared with its base model on the same items. We report the paired difference with a 95\% bootstrap interval (5,000 resamples of items) and test it with the exact McNemar test~\cite{mcnemar1947} for accuracy and a two-sided paired sign-flip permutation test (100,000 permutations) for LAVE. The six medical-versus-base tests are the paper's main hypotheses, so we also report Holm-adjusted $p$ values~\cite{holm1979}. Group- and category-level results are descriptive. All numbers in this paper are recomputed from the saved per-item outputs by one script.

%% ===========================================================================
\section{RESULTS AND ANALYSIS}
\label{sec_results}

\subsection{Main result: medical versus base model}
\label{sec_res_main}
Fig.~\ref{fig_pairs} compares each medical model with its base model on all three tasks. Of the six comparisons, one favours the medical model clearly, one points the same way without reaching significance, three are close to zero and one points the other way.

\begin{figure}[tb]
\centering
%% Values: paired bootstrap (5,000 resamples of items, seed 42), exact McNemar
%% (accuracy) and sign-flip permutation test (LAVE), recomputed from the result
%% CSVs by analysis/verify_numbers.js.
\begin{tikzpicture}
\begin{axis}[
  scale only axis, width=3.05cm, height=3.7cm, clip=false,
  xmin=-9, xmax=16, ymin=0.4, ymax=6.6,
  xtick={-5,0,5,10,15}, xticklabel style={font=\scriptsize},
  xlabel={Medical $-$ base (percentage points)}, xlabel style={font=\scriptsize},
  ytick={1.5,3.5,5.5}, yticklabel style={font=\scriptsize, align=right},
  yticklabels={BanglaMedVQA\\(LAVE), RadImageNet-\\VQA (acc.), BUS-BRA\\(accuracy)},
  major tick length=2pt, axis line style={black!60},
  legend style={font=\scriptsize, at={(0.5,1.03)}, anchor=south, legend columns=2, draw=none,
                /tikz/every even column/.append style={column sep=6pt}},
]
\draw[black!20] (axis cs:-9,2.5) -- (axis cs:16,2.5);
\draw[black!20] (axis cs:-9,4.5) -- (axis cs:16,4.5);
\draw[dashed, black!55] (axis cs:0,0.4) -- (axis cs:0,6.6);
\addplot[only marks, mark=*, mark size=1.8pt, color=cMedVLM,
  error bars/.cd, x dir=both, x explicit, error bar style={line width=0.8pt}]
  coordinates {(0,6) += (0,0) -= (0,0)  (0.8,4) += (2.0,0) -= (2.0,0)  (0.84,2) += (1.16,0) -= (1.14,0)};
\addplot[only marks, mark=square*, mark size=1.7pt, color=cMedGemma,
  error bars/.cd, x dir=both, x explicit, error bar style={line width=0.8pt}]
  coordinates {(7.25,5) += (7.75,0) -= (7.55,0)  (-1.8,3) += (5.4,0) -= (5.4,0)  (7.38,1) += (2.62,0) -= (2.58,0)};
\legend{\medvlm{} $-$ \qwen{}, \medgemma{} $-$ \gemma{}}
\begin{scope}[every node/.style={anchor=west, font=\tiny, align=left, inner sep=0pt}]
\node at (axis cs:16.6,6) {$0.0$; identical on all 400};
\node at (axis cs:16.6,5) {$+7.3$ ($-0.3$, $15.0$)\\ $p=0.08$; 35\% identical};
\node at (axis cs:16.6,4) {$+0.8$ ($-1.2$, $2.8$)\\ $p=0.57$; 93\% identical};
\node at (axis cs:16.6,3) {$-1.8$ ($-7.2$, $3.6$)\\ $p=0.56$; 55\% identical};
\node at (axis cs:16.6,2) {$+0.8$ ($-0.3$, $2.0$)\\ $p=0.16$; 72\% identical};
\node at (axis cs:16.6,1) {$\mathbf{+7.4}$ ($4.8$, $10.0$)\\ $p<0.001$; 36\% identical};
\end{scope}
\end{axis}
\end{tikzpicture}
\caption{Medical model minus its base model on the same items: difference in accuracy (BUS-BRA, RadImageNet-VQA) or mean LAVE (BanglaMedVQA), in percentage points, with 95\% paired bootstrap interval (horizontal bars; values in parentheses on the right); unadjusted $p$ (exact McNemar for accuracy, paired permutation test for LAVE); share of items with identical answers. After Holm correction over the six comparisons only \medgemma{} $-$ \gemma{} on BanglaMedVQA is significant (adjusted $p<0.001$; the next smallest is 0.41).}
\label{fig_pairs}
\end{figure}

\textbf{\medvlm{} behaved almost like a copy of \qwen{}.} It gave identical answers on every BUS-BRA image, on 93\% of RadImageNet-VQA questions and on 72\% of BanglaMedVQA questions, and none of its differences from \qwen{} is distinguishable from zero. A plausible reason is the small scale of its adaptation: 600 MRI multiple-choice questions, with a reward for writing a reasoning block before a letter answer. Our prompts asked for a bare answer, the model complied, and in this direct-answer mode its answers stayed close to its base model's. Whether its reasoning mode would behave differently we did not test (Section~\ref{sec_limitations}).

\textbf{\medgemma{} changed its answers a lot, but not consistently for the better.} Only 35--55\% of its answers matched \gemma{}'s, so its medical training changed behaviour far more. The changes helped clearly in Bangla, possibly on BUS-BRA, and not on RadImageNet-VQA. The share of identical answers is a cheap and useful diagnostic here: it separates an adaptation that left the model's direct answers unchanged from one that changed them without improving them. The next three subsections show what each task adds.

\subsection{BUS-BRA: accuracy hides constant answers}
\label{sec_res_busbra}
Three of the four models scored at chance on BUS-BRA (Table~\ref{tab_busbra}), and the reason is visible in their answers: \qwen{} and \medvlm{} answered ``benign'' for all 400 images, and \gemma{} answered ``malignant'' for 391. On a balanced sample a constant answer earns exactly 50\%. \gemma{}'s 98\% sensitivity therefore means nothing on its own: with 2.5\% specificity it flags almost every patient.

\medgemma{} was the only model whose answers varied substantially across images. It reached 57.5\% (95\% CI 52.6--62.3\%), significantly above chance ($p=0.0016$, still significant after Bonferroni correction for four models). It leaned towards ``benign'', however, and missed 118 of the 200 cancers. Its gain over \gemma{} (+29 images) was not significant in the paired test (144 images correct only for \medgemma{} against 115 only for \gemma{}). So \medgemma{} shows a weak signal that its base model lacks, but this sample does not establish that it is better than \gemma{}, nor that its medical training is the reason.

\begin{table}[tb]
\centering
\caption{BUS-BRA results (400 images, 200 per class). ``Malig.'' = number of ``malignant'' answers; Sens./Spec.\ = sensitivity/specificity for malignant (\%); $p$: one-sided binomial test against 50\%.}
\label{tab_busbra}
\footnotesize
\setlength{\tabcolsep}{2.5pt}
\resizebox{\columnwidth}{!}{%
\begin{tabular}{@{}lccccc@{}}
\toprule
\textbf{Model} & \textbf{Acc.\ \% (95\% CI)} & \textbf{Malig.} & \textbf{Sens.} & \textbf{Spec.} & $p$ \\
\midrule
\qwen{}-2B & 50.0 (45.1--54.9) & 0 & 0.0 & 100.0 & 0.52 \\
\medvlm{} & 50.0 (45.1--54.9) & 0 & 0.0 & 100.0 & 0.52 \\
\gemma{}~4B & 50.3 (45.4--55.1) & 391 & 98.0 & 2.5 & 0.48 \\
\medgemma{}~4B & \textbf{57.5} (52.6--62.3) & 134 & 41.0 & 74.0 & 0.002 \\
\bottomrule
\end{tabular}}
\end{table}

\subsection{RadImageNet-VQA: anatomy is easy, disease is not}
\label{sec_res_radimagenet}
Overall accuracy was similar for all four models: 64.2\% (\qwen{}), 65.0\% (\medvlm{}), 63.0\% (\gemma{}) and 61.2\% (\medgemma{}), against a chance level of 40\%. The question groups tell a clearer story (Fig.~\ref{fig_radimagenet}). Every model recognised anatomy well (79--88\% on both anatomy groups), but on questions about disease all of them were near chance: 52--61\% on abnormality and pathology yes/no (chance 50\%) and 32--37\% on pathology multiple choice (chance 25\%). With 100 questions per group the 95\% interval is about $\pm$9 points, so several of the disease scores cannot be separated from guessing. Neither medical model narrowed this gap; \medgemma{} was in fact below \gemma{} on both anatomy groups, although neither difference is significant ($p=0.39$ and $0.79$).

\begin{figure}[tb]
\centering
\begin{tikzpicture}
\begin{axis}[modelbars, ymax=100, ylabel={Accuracy (\%)},
  symbolic x coords={AM,AY,BY,PY,PM},
  xticklabels={Anatomy\\MCQ, Anatomy\\yes/no, Abnorm.\\yes/no, Pathology\\yes/no, Pathology\\MCQ}]
\addplot[fill=cQwen, draw=black] coordinates {(AM,87) (AY,88) (BY,55) (PY,57) (PM,34)};
\addplot[fill=cMedVLM, draw=black] coordinates {(AM,88) (AY,86) (BY,61) (PY,57) (PM,33)};
\addplot[fill=cGemma, draw=black, gemmapair] coordinates {(AM,84) (AY,81) (BY,52) (PY,61) (PM,37)};
\addplot[fill=cMedGemma, draw=black, gemmapair] coordinates {(AM,80) (AY,79) (BY,55) (PY,60) (PM,32)};
\legend{Qwen2-VL, MedVLM-R1, Gemma 3, MedGemma}
\draw[dashed, gray] ([xshift=-14pt]axis cs:AM,25) -- ([xshift=14pt]axis cs:AM,25);
\draw[dashed, gray] ([xshift=-14pt]axis cs:AY,50) -- ([xshift=14pt]axis cs:AY,50);
\draw[dashed, gray] ([xshift=-14pt]axis cs:BY,50) -- ([xshift=14pt]axis cs:BY,50);
\draw[dashed, gray] ([xshift=-14pt]axis cs:PY,50) -- ([xshift=14pt]axis cs:PY,50);
\draw[dashed, gray] ([xshift=-14pt]axis cs:PM,25) -- ([xshift=14pt]axis cs:PM,25);
\end{axis}
\end{tikzpicture}
\caption{RadImageNet-VQA accuracy by question group (100 questions each). Dashed lines mark chance. Light bars are general models, dark bars their medical versions; the Gemma pair is hatched.}
\label{fig_radimagenet}
\end{figure}

\textbf{When the image does not decide, the models fall back on preferences} (Table~\ref{tab_prefs}). Because each yes/no group was balanced, an unbiased model would say ``yes'' about half the time. \gemma{} said ``yes'' to 74\% of yes/no questions, including 98 of the 100 abnormality questions. \medgemma{}, derived from it, leaned the other way, with 42\% ``yes'' overall and 32\% on pathology. The medical version did not lose the preference but showed the opposite one, and on abnormality questions both models ended up near chance from opposite directions. On pathology multiple choice the correct letters were spread evenly (A~25, B~22, C~25, D~28), yet each pair favoured one letter: B for the Gemma pair, C for the Qwen pair. This matches the option-letter bias known for language models~\cite{zheng2024mcqbias} and suggests the preference is inherited from the base model.

\begin{table}[tb]
\centering
\caption{RadImageNet-VQA answer preferences. Share of ``yes'' answers per yes/no group (true share 50\%), and the most chosen letter on pathology MCQ with its count (the correct letter was B in 22 of 100 questions, C in 25).}
\label{tab_prefs}
\footnotesize
\setlength{\tabcolsep}{3pt}
\resizebox{\columnwidth}{!}{%
\begin{tabular}{@{}lcccc@{}}
\toprule
 & \textbf{Qwen2} & \textbf{MedVLM} & \textbf{Gemma3} & \textbf{MedGemma} \\
\midrule
``Yes'', anatomy (\%) & 42 & 42 & 49 & 55 \\
``Yes'', abnormality (\%) & 77 & 69 & 98 & 39 \\
``Yes'', pathology (\%) & 43 & 47 & 75 & 32 \\
``Yes'', all yes/no (\%) & 54 & 53 & 74 & 42 \\
\midrule
Top letter, path.\ MCQ & C (37) & C (35) & B (35) & B (42) \\
\bottomrule
\end{tabular}}
\end{table}

The four models also failed together: all of them were wrong on the same 65 questions (13\%). Accuracy varied by body region (hip was hardest for every model, 37--53\% on 38 questions), but most region-level cells are too small to interpret.

\subsection{BanglaMedVQA: the one clear gain, and why it needs care}
\label{sec_res_bangla}
BanglaMedVQA is the only task where a medical model clearly beat its base model (Table~\ref{tab_bangla}): \medgemma{} reached 31.5\% LAVE against 24.1\% for \gemma{} (Holm-adjusted $p<0.001$). The two Qwen-family models were close to each other at about 15--16\%. Exact match was near zero for all models, which is why LAVE carries the comparison, and the rest of this subsection asks how far LAVE can be trusted.

\begin{table}[tb]
\centering
\caption{BanglaMedVQA results (500 questions). EM = exact match (\%). LAVE with 95\% bootstrap interval. ``No Bangla'' = share of answers containing no Bangla script (\%).}
\label{tab_bangla}
\footnotesize
\setlength{\tabcolsep}{3pt}
\resizebox{\columnwidth}{!}{%
\begin{tabular}{@{}lcccc@{}}
\toprule
\textbf{Model} & \textbf{EM} & \textbf{BERTScore} & \textbf{LAVE (95\% CI)} & \textbf{No Bangla} \\
\midrule
\qwen{}-2B & 0.0 & 68.6 & 15.1 (12.5--17.8) & 82.0 \\
\medvlm{} & 0.0 & 68.4 & 16.0 (13.4--18.7) & 89.2 \\
\gemma{}~4B & 2.4 & 69.0 & 24.1 (20.9--27.3) & 90.0 \\
\medgemma{}~4B & 1.6 & \textbf{71.4} & \textbf{31.5} (28.1--34.9) & 98.2 \\
\bottomrule
\end{tabular}}
\end{table}

\textbf{The gain is in the general categories.} By category, BanglaMedVQA repeats the anatomy-versus-disease pattern of RadImageNet-VQA. Modality was the easiest category for every model (LAVE 55.8--73.3), and organ questions were answered well only by the Gemma pair (47.1 for \gemma{}, 63.1 for \medgemma{}, against 3.3--5.3 for the Qwen pair). Condition, abnormality and position stayed between 0 and 13 for all four models in both sources, although the abnormality scores mostly reflect the scoring problem described below. \medgemma{}'s lead over \gemma{} came almost entirely from chest X-ray organ (+30.6 points) and chest X-ray modality (+20.4) questions, the only two source--category cells whose paired intervals clearly exclude zero. On organ questions the difference is mostly real: \medgemma{} answered ``Chest'' to 48 of 50, while \gemma{} often named the lung instead of the chest given in the reference, and three times answered with the Chinese characters for ``chest''. The modality difference is not: by the rule-based check both models named the modality correctly on almost every chest X-ray (\gemma{} 50, \medgemma{} 49 of 50), and the gap comes from how the judge scored \gemma{}'s answers (below). Either way, the one clear gain concerns general questions and says little about diagnostic ability.


\textbf{The models mostly answered in English.} Although every question was in Bangla, 82--98\% of each model's answers contained no Bangla script; \medgemma{}, the best model, answered in Bangla least often. When the Qwen models did answer in Bangla the answer was often meaningless; their most frequent answer to chest X-ray organ questions was \bn{লাইন} (``line''). They also answered some MedICaT questions with a bare option letter although no options were given (39 answers for \qwen{}, 27 for \medvlm{}). Several models repeated one stock answer: \qwen{} and \medvlm{} each said ``pneumonia'' to 37 of the 50 chest X-ray condition questions. On the balanced abnormality questions, the yes/no preferences from RadImageNet-VQA appeared again: the Qwen models never answered ``yes'', while 62 of \medgemma{}'s 79 English yes/no answers were ``Yes''.

\textbf{The metric under-credits correct English answers} (Table~\ref{tab_audit}). The LAVE judge sees only the Bangla reference and the model's answer, so it has to decide on its own whether an English answer matches a Bangla reference. \medgemma{} answered 34 Bangla ``yes'' questions with the English ``Yes''; every one of these correct answers scored 0, while its five correct answers in Bangla (three \bn{হ্যাঁ}, ``yes'', and two \bn{না}, ``no'') all scored 1.0. The judge was also inconsistent across models: some of the Qwen models' correct English ``no'' answers received credit, but none of the Gemma models' correct English yes/no answers did. On modality questions, correct ``CT'' answers averaged 0.26--0.33 LAVE, against 0.80--0.87 for correct ``X-ray'' and 0.81--0.86 for correct ``MRI'' answers, for every model; \gemma{}'s 14 answers of ``Radiograph(y)'' to X-ray questions averaged 0.04. BERTScore without rescaling did not help either: answers with zero LAVE still averaged 66--70, and its rank correlation with LAVE was only 0.27--0.43.

\begin{table}[tb]
\centering
\caption{Audit of LAVE on answers that are correct by a rule-based check. Top: correct English yes/no answers to abnormality questions (an answer starting with ``yes'' or ``no'', e.g.\ ``No abnormality'', or ``There is no abnormality'') and how many the judge credited (LAVE $>0$); \medgemma{}'s five correct Bangla yes/no answers, all credited, are not included. Bottom: correct modality answers (of 100) by keyword matching of X-ray/radiograph, CT, MRI or ultrasound in either language.}
\label{tab_audit}
\footnotesize
\setlength{\tabcolsep}{3pt}
\resizebox{\columnwidth}{!}{%
\begin{tabular}{@{}lcccc@{}}
\toprule
 & \textbf{Qwen2} & \textbf{MedVLM} & \textbf{Gemma3} & \textbf{MedGemma} \\
\midrule
Correct yes/no answers & 31 & 34 & 12 & 58 \\
\quad credited by LAVE & 8 & 11 & 0 & 5 \\
\midrule
Correct modality (rule) & 82 & 84 & 99 & 99 \\
Modality LAVE (\%) & 55.8 & 56.5 & 60.7 & 73.3 \\
\bottomrule
\end{tabular}}
\end{table}

So the LAVE scores in Table~\ref{tab_bangla} are noisy and under-credit correct answers given in English, and the size of the gap between models partly depends on which language they answer in. This matches the poor cross-lingual consistency reported for LLM judges in general~\cite{fu2025multilingualjudge}. It also changes the size of the one clear gain. If modality answers are scored by the rule-based check instead of LAVE, \medgemma{}'s lead over \gemma{} shrinks from 7.4 to 4.9 points (95\% CI 2.4--7.3; $p<0.001$), and the \medvlm{}--\qwen{} difference stays non-significant (+1.1, $p=0.11$). The advantage survives, but about a third of it was a scoring artefact. The scoring problem also limits comparison with the original paper, which reported 19.85\% zero-shot LAVE for \medgemma{}~4B on the test split~\cite{ahmed2026banglamedvqa}, against our 31.5\% on a different sample.

%% ===== REVISION (blue): blank-image control results, Reviewers 9 and 11 ======
\subsection{\textcolor{blue}{Blank-image control: how much do the answers depend on the image?}}
\label{sec_res_blank}
{\color{blue}
Table~\ref{tab_blank} compares each model's score on the real images with its score on the blank image. On RadImageNet-VQA and BanglaMedVQA every model scored much lower with the blank image (RadImageNet-VQA accuracy: McNemar $p<10^{-10}$ for all four models; BanglaMedVQA LAVE: paired permutation $p<0.001$ for all four). On BUS-BRA the blank image changed the accuracy only for \medgemma{}.
}

\begin{table}[tb]
\color{blue}
\centering
\caption{Real-image $\rightarrow$ blank-image scores. Accuracy (\%) on BUS-BRA (chance 50\%) and RadImageNet-VQA (chance 40\% over the 500 questions), and mean LAVE (\%) on BanglaMedVQA. The blank image is white for RadImageNet-VQA and BanglaMedVQA; on BUS-BRA the white and black images gave identical results.}
\label{tab_blank}
\footnotesize
\setlength{\tabcolsep}{3pt}
\resizebox{\columnwidth}{!}{%
\begin{tabular}{@{}lccc@{}}
\toprule
\textbf{Model} & \textbf{BUS-BRA acc.} & \textbf{RadImageNet acc.} & \textbf{Bangla LAVE} \\
\midrule
\qwen{}-2B & 50.0 $\rightarrow$ 50.0 & 64.2 $\rightarrow$ 41.6 & 15.1 $\rightarrow$ 2.7 \\
\medvlm{} & 50.0 $\rightarrow$ 50.0 & 65.0 $\rightarrow$ 43.8 & 16.0 $\rightarrow$ 2.6 \\
\gemma{}~4B & 50.3 $\rightarrow$ 50.0 & 63.0 $\rightarrow$ 38.8 & 24.1 $\rightarrow$ 11.7 \\
\medgemma{}~4B & 57.5 $\rightarrow$ 50.0 & 61.2 $\rightarrow$ 38.4 & 31.5 $\rightarrow$ 7.2 \\
\bottomrule
\end{tabular}}
\end{table}

{\color{blue}
\textbf{RadImageNet-VQA: anatomy depends on the image, disease barely does.} With the blank image, overall accuracy fell to 38.4--43.8\%, and the 95\% interval of every model includes the 40\% chance level. The loss was concentrated in the anatomy groups (Table~\ref{tab_blank_groups}): anatomy MCQ fell from 80--88\% to 11--36\% and anatomy yes/no from 79--88\% to 51--57\%, a gain from seeing the image of 28 to 73 points. Disease questions were already close to chance with the real image and moved little: pathology MCQ was 24--25\% (chance 25\%) against 32--37\%, and abnormality and pathology yes/no were 50--54\% against 52--61\%. The gain from the image on disease questions was between $-2$ and $+13$ points for every model and group. Two blank-image anatomy MCQ scores lay below the 25\% chance level (\gemma{} 11\%, \medgemma{} 15\%) because these models favoured letter B (63 and 77 of 100 answers), which is correct for only 22 of the 100 questions.
}

\begin{table}[tb]
\color{blue}
\centering
\caption{RadImageNet-VQA accuracy (\%) by group, real $\rightarrow$ blank image (100 questions per group; chance 25\% for MCQ and 50\% for yes/no).}
\label{tab_blank_groups}
\footnotesize
\setlength{\tabcolsep}{3pt}
\resizebox{\columnwidth}{!}{%
\begin{tabular}{@{}lcccc@{}}
\toprule
 & \textbf{Qwen2} & \textbf{MedVLM} & \textbf{Gemma3} & \textbf{MedGemma} \\
\midrule
Anatomy, MCQ & 87 $\rightarrow$ 28 & 88 $\rightarrow$ 36 & 84 $\rightarrow$ 11 & 80 $\rightarrow$ 15 \\
Anatomy, yes/no & 88 $\rightarrow$ 55 & 86 $\rightarrow$ 57 & 81 $\rightarrow$ 52 & 79 $\rightarrow$ 51 \\
Abnormality, yes/no & 55 $\rightarrow$ 50 & 61 $\rightarrow$ 50 & 52 $\rightarrow$ 54 & 55 $\rightarrow$ 50 \\
Pathology, yes/no & 57 $\rightarrow$ 51 & 57 $\rightarrow$ 51 & 61 $\rightarrow$ 53 & 60 $\rightarrow$ 51 \\
Pathology, MCQ & 34 $\rightarrow$ 24 & 33 $\rightarrow$ 25 & 37 $\rightarrow$ 24 & 32 $\rightarrow$ 25 \\
\bottomrule
\end{tabular}}
\end{table}

{\color{blue}
\textbf{Which answer preferences come from the image?} The control separates preferences that depend on the input from those that do not. The option-letter preference does not: on pathology MCQ the most chosen letter with the blank image was the same as with the real image for all four models (C for the Qwen pair, B for the Gemma pair), and \gemma{} chose B 54 times instead of 35, so this preference behaves like a language prior inherited from the base model. The yes/no preferences do depend on the image. The share of ``yes'' answers over the 300 yes/no questions fell from 54\%, 53\%, 74\% and 42\% to 12\%, 13\%, 30\% and 14\% (\qwen{}, \medvlm{}, \gemma{}, \medgemma{}), and on abnormality questions no model said ``yes'' more than 4 times in 100 (real image: 77, 69, 98 and 39). \gemma{}'s habit of saying ``yes'' to almost every abnormality question therefore appears only when a scan is shown; with a blank image all four models mostly answered ``no''.

BUS-BRA shows the same pattern. With a blank image, white or black and in every repetition, all four models answered ``benign'' for all 400 cases. \qwen{} and \medvlm{} gave this answer for the real images as well (identical on all 400 cases), so their BUS-BRA output does not depend on the image at all. \gemma{} answered ``malignant'' for 391 real images but ``benign'' for the blank one (2\% identical answers), so its ``malignant'' preference is triggered by the ultrasound image rather than by the prompt. \medgemma{} gave the same answer as with the blank image on 66\% of cases, and its 57.5\% accuracy on real images exceeds the 50.0\% of the blank image (82 cases correct only with the real image, 52 only with the blank one; exact McNemar $p=0.012$, unadjusted).

\textbf{BanglaMedVQA: \medgemma{} depends on the image more than \gemma{}.} Mean LAVE fell by 12.4, 13.4, 12.4 and 24.3 points for \qwen{}, \medvlm{}, \gemma{} and \medgemma{} (95\% bootstrap intervals 9.7--15.3, 10.8--16.2, 9.0--15.8 and 20.7--28.0), and exact match was 0 of 500 for every model except \gemma{} (2 of 500). The blank-image answers were nearly constant within a category. \gemma{} answered ``X-ray'' to 94 of the 100 modality questions, \medgemma{} answered ``CT'' to all 100 and ``brain'' to 98 of the 100 organ questions, and \qwen{} and \medvlm{} named no modality at all (0 of 100 correct by the rule-based check of Section~\ref{sec_res_bangla}; typical answers ``this is a medical image'' or ``digital''); \qwen{} answered \bn{লাইন} (``line'') to 33 of the 100 organ questions. A constant answer earns the share of references it happens to match: X-ray is correct for 56\% of the modality references, so \gemma{} kept a modality LAVE of 48.7 on the blank image (real image 60.7), whereas \medgemma{} fell from 73.3 to 7.3 and \qwen{} and \medvlm{} from 55.8 and 56.5 to 1.6 and 2.3. Organ LAVE fell from 47.1 to 0.5 for \gemma{} and from 63.1 to 21.4 for \medgemma{}. On the three specialised categories (abnormality, condition, position) LAVE was 0.0--11.1 with the real images and 0.3--7.3 with the blank image. \gemma{} also returned an empty answer for 125 of the 500 blank-image questions (11 with real images), which lowers its blank-image scores.

The real-minus-blank LAVE of \medgemma{} is 11.9 points larger than that of \gemma{} (95\% CI 7.2--16.5; paired permutation $p<0.001$), whereas \medvlm{} and \qwen{} differ by only 0.9 points ($-0.6$ to $2.4$; $p=0.23$). With the blank image \medgemma{} scored lower than \gemma{} (7.2 against 11.7), so its 7.4-point advantage with real images is not a language prior: it appears only when an image is shown.

\textbf{What the control changes.} It leaves the paired comparison unchanged: \medvlm{} remains indistinguishable from \qwen{}, and the \medgemma{} advantage in Bangla is image-dependent. It sharpens two statements. Anatomy and modality answers depend on the image, whereas answers to disease questions were only slightly better with the real image than with a blank one. And the answer preferences are of two kinds: the option-letter preference is independent of the image, while the ``yes'' and ``malignant'' preferences appear only when a scan is shown. A blank image is an unnatural input and we used one white image per task (white and black on BUS-BRA), so the control does not replace a shuffled-image test.
}
%% ===== end of revision ======================================================


\subsection{Discussion}
\textbf{What the three tasks show together.} The four models behaved alike in two ways. They recognised \emph{what kind of image} they were given (modality, anatomy) much better than \emph{what was wrong} in it, and when the image did not settle the answer they used a fixed preference: one class label, ``yes'' or ``no'', one option letter or one stock disease name. Each medical model either shared its base model's preference or showed the opposite one; none lost it. Against this background, the medical versions changed little. On BUS-BRA, \medvlm{} gave exactly its base model's answers and \medgemma{}'s improvement over a chance-level base was not significant. On RadImageNet-VQA, neither closed the gap between anatomy and disease, even though \medvlm{} was trained on MRI questions of a similar style. BanglaMedVQA gave the only clear gain, on general rather than diagnostic questions, and part of it came from the metric.

This agrees with the controlled comparison of earlier medical VLMs by Jeong et al.~\cite{jeong2025limited}, obtained there with model-specific prompts and here with a shared prompt, and extends it to a reinforcement-learning adaptation, a large multimodal adaptation and a low-resource language. It does not contradict the gains reported by the \medgemma{} authors~\cite{sellergren2025medgemma}, which were measured on other benchmarks with their own prompts; it shows that such gains do not transfer automatically to every medical task or prompt.

\textbf{Lessons for evaluating medical VLMs.} Three practices made the difference between a misleading and an informative evaluation, and each is cheap. First, compare a medical model with its own base model under identical conditions; without that baseline, \medgemma{}'s 57.5\% on BUS-BRA or 31.5\% LAVE in Bangla would look like the effect of medical training. Second, use balanced labels and report the distribution of answers: on BUS-BRA's natural class mix (68\% benign) the always-``benign'' Qwen models would have scored about 68\%, and on any yes/no set skewed towards ``yes'', \gemma{}'s ``yes'' preference would look like skill. The share of identical answers within a pair is similarly cheap and showed at once that one adaptation had left the model's direct answers unchanged. Third, audit LLM-judged metrics whenever models may answer in a different language from the reference; a rule-based check on a few closed question types was enough to show that the judge's verdicts depended on answer language and wording.

For the two small models we tested, the medical training their authors applied did not translate into better zero-shot answers to diagnostic questions. Our design cannot say which part of that training would need to change, but the results are consistent with a need for training methods and data aimed specifically at diagnostic questions and at low-resource languages.

%% ===========================================================================
\section{LIMITATIONS}
\label{sec_limitations}

\textbf{One shared, direct-answer prompt per task.} An identical prompt isolates the effect of the weights, but does not give each model its best conditions. This matters most for \medvlm{}, which was trained to reason in \texttt{<think>} tags before answering (its model card generates up to 1,024 tokens). We tested it only in direct-answer mode, so our results do not show whether its reasoning mode helps; they show that its direct answers are almost unchanged. \medgemma{} might also do better with its own recommended prompts.

% \textbf{What the pairs compare.} Each medical model differs from its base model in several ways at once (training data, objective and, for \medgemma{}, a medically tuned vision encoder). Our comparison measures the released models, not the effect of any single ingredient, and two pairs are too few to generalise to medical VLMs as a class.

\textbf{Small samples and one run.} Each RadImageNet-VQA group has 100 questions and each BanglaMedVQA cell 50, so group-level intervals are wide. Each item was run once with greedy decoding.

% \textbf{Text outputs only, no image-ablation control.} We stored decoded answers, not probabilities, so we cannot compute AUC on BUS-BRA; a model with a strong label preference might still rank lesions well. We also did not run blank-image controls, which would show how much each answer depends on the image.

\textbf{Bangla evaluation.} LAVE relies on a proprietary judge that sees no question and scored correct English answers inconsistently, and BERTScore was computed without rescaling. The BanglaMedVQA references were LLM-generated and machine-translated, with expert checks on a subset only. We did not run the English versions of the questions, which would separate visual errors from language errors, and our train-split sample differs from the original paper's test split.

% \textbf{Possible training-data overlap.} \medvlm{}'s training questions come from the OmniMedVQA collection; example questions on its model card use labels such as ``flexor pathology'' that also appear in our RadImageNet-VQA sample. \medgemma{}'s training data are not fully public. Overlap, if any, would favour the medical models.

\textbf{Reproducibility and scope.} Model revisions and library versions were pinned for BUS-BRA but not recorded for the two VQA tasks. We tested two pairs of 2B and 4B models on three datasets; larger models or fine-tuned use may behave differently. None of these results is a clinical evaluation.

%% ===========================================================================
\section{CONCLUSIONS}
\label{sec_conclusion}
We compared two small medical VLMs with their own base models on breast ultrasound, CT/MRI and Bangla medical VQA, using identical items, prompts and greedy decoding. In this zero-shot, direct-answer setting, medical adaptation did not consistently improve performance. MedVLM-R1 behaved almost like Qwen2-VL-2B. MedGemma 4B was clearly better than Gemma 3 4B only on general Bangla questions, and about a third of that gain came from how the LLM judge scored English answers. Our findings cover two model pairs, one shared prompt and one run, and do not test MedVLM-R1’s reasoning mode, so they do not show that medical training is ineffective. {\color{blue}Blank-image controls showed that the models' anatomy and modality answers depend on the image, whereas answers to disease questions were only slightly better with the real image than with a blank one. Future work should add per-model prompt tuning, shuffled-image controls, English versions of the Bangla questions and larger models.}




%% Double-blind: no acknowledgements in the review version.

\bibliographystyle{ijcvsplatex-num}
\bibliography{references}

\end{document}
